Segons el model atòmic de Bohr, les òrbites que descriuen els electrons al voltant del nucli compleixen la relació següent:
\[ m_\mathrm{e} v r = n\,\frac{h}{2\pi}, \]
on $m_\mathrm{e}$ és la massa de l'electró, $v$ la seva velocitat, $r$ el radi de l'òrbita i $n$ és un nombre enter (nombre quàntic principal). Per a l'estat fonamental, $n$ val 1; per a la segona òrbita, $n$ val 2, i així successivament.

\figura[0.3]{fig1.png}

\begin{apartat}{1,25}
En l'hidrogen muònic, el muó substitueix l'electró. El muó és una partícula idèntica a l'electró (és un leptó com l'electró i té la mateixa càrrega), però la seva massa és unes 200 vegades la massa de l'electró. Determineu el radi de l'òrbita del muó en el seu estat fonamental ($n = 1$). Quin àtom ocupa un volum més gran, l'àtom d'hidrogen o l'hidrogen muònic? Justifiqueu la resposta.
\end{apartat}

\begin{apartat}{1,25}
L'energia del muó en l'estat fonamental és $-4,355\times 10^{-16}\ \mathrm{J}$ i en la segona òrbita ($n = 2$) és $-1,089\times 10^{-16}\ \mathrm{J}$. Quan el muó passa de la segona òrbita a l'estat fonamental emet un fotó. Quina és la freqüència i la longitud d'ona d'aquest fotó?
\end{apartat}

\dades{%
  $m_\mathrm{e} = 9,11\times 10^{-31}\ \mathrm{kg}$.\\
  $|e| = 1,602\times 10^{-19}\ \mathrm{C}$.\\
  $h = 6,63\times 10^{-34}\ \mathrm{J\,s}$.\\
  $c = 3,00\times 10^{8}\ \mathrm{m\,s^{-1}}$.\\
  $k = \dfrac{1}{4\pi\varepsilon_0} = 8,99\times 10^{9}\ \mathrm{N\,m^2\,C^{-2}}$.}
